\newpage
\begin{tcolorbox}[title=Problem 1, breakable]
\begin{enumerate}
    \item If $P$ and $Q$ are distinct rational points in the $(x, y)$ plane,
    prove that the line connecting them is a rational line.
    \item If $L_1$ and $L_2$ are distinct rational lines in $(x, y)$ plane,
    prove that their intersection is a rational point (or empty).
\end{enumerate}
\end{tcolorbox}

\begin{proof}
    Suppose $P$ and $Q$ are distinct rational points in the $(x, y)$ plane.
    Let $P = (p_1, p_2)$ and $Q = (q_1, q_2)$.
    We seek an equation of a line with rational coefficients
        such that $P$ and $Q$ are solutions.
    There are two cases to consider.
    First, suppose $q_1 = p_1$.
    Since $P \ne Q$ we must have $p_2 \ne q_2$.
    Then the vertical line $x = p_1$ has rational coefficients
        and is satisfied by both $P$ and $Q$, since $p_1 = q_1$ is rational.
    Now, suppose $q_1 \ne p_1$.
    Set $m = \frac{q_2 - p_2}{q_1 - p_1}$ and $b = p_2 - m p_1$.
    Both are rational since $p_1, p_2, q_1, q_2$ are rational and $q_1 - p_1 \ne 0$.
    Substituting the coordinates of $P$ into $y = mx + b$ gives
    \[p_2 = \left( \frac{q_2 - p_2}{q_1 - p_1} \right) p_1 + b,\]
    which holds by our choice of $b$. Substituting $Q$ gives
    \[m q_1 + b = m q_1 + p_2 - m p_1 = m(q_1 - p_1) + p_2
        = (q_2 - p_2) + p_2 = q_2,\]
    so $Q$ is also a solution.
\end{proof}

\begin{proof}
    Suppose $L_1$ and $L_2$ are distinct rational lines in the $(x, y)$ plane.
    Write $L_1$ and $L_2$ as
    \[y = m_1 x + b_1 \text{ and } y = m_2 x + b_2 \text{ respectively,}\]
    where $m_1, m_2, b_1, b_2$ are rational.
    A point $(x, y)$ lies on both lines exactly when it solves the system
    \[
    \begin{bmatrix}
        -m_1 & 1 \\
        -m_2 & 1
    \end{bmatrix}
    \begin{bmatrix}
        x \\ y
    \end{bmatrix} =
    \begin{bmatrix}
        b_1 \\ b_2
    \end{bmatrix}.
    \]
    Let $A = \begin{bmatrix} -m_1 & 1 \\ -m_2 & 1 \end{bmatrix}$, so $\det A = -m_1 + m_2 = m_2 - m_1$.
    Suppose $m_1 \ne m_2$, so $\det A \ne 0$ and $A$ is invertible.
    By Cramer's rule the system has the unique solution
    \[
    x = \frac{1}{\det A}\begin{vmatrix} b_1 & 1 \\ b_2 & 1 \end{vmatrix}
        = \frac{b_1 - b_2}{m_2 - m_1},
    \qquad
    y = \frac{1}{\det A}\begin{vmatrix} -m_1 & b_1 \\ -m_2 & b_2 \end{vmatrix}
        = \frac{m_2 b_1 - m_1 b_2}{m_2 - m_1}.
    \]
    Since $m_1, m_2, b_1, b_2$ are rational and $m_2 - m_1 \ne 0$, both $x$ and $y$ are rational,
    so the intersection point is rational.
\end{proof}

\newpage
\begin{tcolorbox}[title=Problem 2, breakable]
    Let $C$ be the conic section given by the equation 
    \[F(x, y) = ax^2 + bxy + cy^2 + dx + ey + f = 0,\]
    and let $\delta$ be the determinant 
    \[det \begin{bmatrix}
        2a & b & d \\
        b & 2c & e \\
        d & e & 2f
    \end{bmatrix}.\]
    \begin{enumerate}
        \item Show that if $\delta \ne 0$, then $C$ has no singular points. That is shows that there 
        are no points $(x, y)$ where 
        \[F(x, y) = \frac{\partial F}{\partial x}(x, y) = \frac{\partial F}{\partial y} = 0,\]
        \item Conversely, show that if $\delta = 0$ and $b^2 - 4ac \ne 0$, then there is a unique singular point on $C$.
        \item Let $L$ be the line $y = \alpha x + \beta$ with $\alpha \ne 0$. Show that the intersection 
        of $L$ and $C$ consists of either zero, one, or two points.
        \item Determine the conditions on the coefficients to ensure that the intersection $L \cap C$
        consists of exactly one point. What is the geometric significance of these conditions?
        (Note. That there will be more than one case.)
    \end{enumerate}
\end{tcolorbox}

\begin{proof}
    Let $(x, y)$ be a singular point of $C$, so that $F = 0$, $\frac{\partial F}{\partial x} = 0$,
    and $\frac{\partial F}{\partial y} = 0$. Since
    \[\frac{\partial F}{\partial x} = 2ax + by + d \text{ and } \frac{\partial F}{\partial y} = bx + 2cy + e,\]
    and since
    \[
    2F =
    \begin{bmatrix}
        x & y & 1
    \end{bmatrix}
    \begin{bmatrix}
        2a & b & d \\
        b & 2c & e \\
        d & e & 2f
    \end{bmatrix}
    \begin{bmatrix}
        x \\ y \\ 1
    \end{bmatrix},
    \]
    the three conditions $\frac{\partial F}{\partial x} = \frac{\partial F}{\partial y} = 0$ and $dx + ey + 2f = 0$ give
    \[
    \begin{bmatrix}
        2a & b & d \\
        b & 2c & e \\
        d & e & 2f
    \end{bmatrix}
    \begin{bmatrix}
        x \\ y \\ 1
    \end{bmatrix}
    =
    \begin{bmatrix}
        0 \\ 0 \\ 0
    \end{bmatrix}.
    \]
    The vector $(x, y, 1) \ne \theta$ so $\delta = 0$ which is a contradiction.
\end{proof}

\begin{proof}
    Suppose $\delta = 0$ and $b^2 - 4ac \ne 0$. As above, a point $(x, y)$ is singular iff
    \[
    \begin{bmatrix}
        2a & b & d \\
        b & 2c & e \\
        d & e & 2f
    \end{bmatrix}
    \begin{bmatrix}
        x \\ y \\ 1
    \end{bmatrix}
    =
    \begin{bmatrix}
        0 \\ 0 \\ 0
    \end{bmatrix}.
    \]
    The first two rows say
    \[
    \begin{bmatrix}
        2a & b \\
        b & 2c
    \end{bmatrix}
    \begin{bmatrix}
        x \\ y
    \end{bmatrix}
    =
    \begin{bmatrix}
        -d \\ -e
    \end{bmatrix},
    \]
    and the determinant of this $2 \times 2$ system is $4ac - b^2 \ne 0$, so it has a unique solution $(x, y)$.
    Since $\delta = 0$ but the first two rows are independent, the third row is a combination of them,
    so it vanishes on $(x, y, 1)$. Thus $(x, y)$ satisfies $dx + ey + 2f = 0$, and is the
    unique singular point on $C$.
\end{proof}

\begin{proof}
    We can write $y = \alpha x + \beta$ since $L$ is given. Then substituting into $C$ we find 
    \[F(x, \alpha x + \beta) 
    = a x^2 
    + b x\left(\alpha x + \beta\right)
    + c\left(\alpha x + \beta\right)^2 
    + d x 
    + e\left(\alpha x + \beta\right)
    + f = 0.\]
    Collecting gives a quadratic in $x$,
    \[
    A x^2 + B x + C = 0,
    \]
    where
    \[A = a + b\alpha + c\alpha^2, \quad B = b\beta + 2c\alpha\beta + d + e\alpha, \quad C = c\beta^2 + e\beta + f.\]
    This clearly has $0, 1$, or $2$ roots.
\end{proof}

\begin{proof}
    With $A, B, C$ as above, $L \cap C$ is one point exactly when $A x^2 + B x + C = 0$ has one solution.
    There are two cases. If $A \ne 0$, the quadratic has one root iff
    \[B^2 - 4AC = 0.\]
    If $A = 0$, that is $a + b\alpha + c\alpha^2 = 0$, the equation is
    linear and has one root iff $B \ne 0$.
\end{proof}

\newpage
\begin{tcolorbox}[title=Problem 3, breakable]
    Let $C$ be the conic given by the equation 
    \[F(x,y) = x^2 - 3xy + 2y^2 - x + 1 = 0.\]
    \begin{enumerate}
        \item Check that $C$ is non-singular (cf. the previous exercise).
        \item Let $L$ be the line $y = \alpha x + \beta$. Suppose that the intersection $L \cap C$ contains the point 
        $P_0 = (x_0, y_0)$. Assuming that $L \cap C$ consists of two distinct points, find the second point 
        of $L \cap C$ in terms of $\alpha, \beta, x_0, y_0$.
        \item If $L$ is a rational line and $P_0$ is a rational point (i.e., $\alpha, \beta, x_0, y_0 \in \mathbb{Q}$),
        prove that the second point of $L \cap C$ is also a rational point.
    \end{enumerate}
\end{tcolorbox}

\begin{proof}
    We have
    \[
    A = \begin{bmatrix}
    1 & -\tfrac{3}{2} & -\tfrac{1}{2} \\[4pt]
    -\tfrac{3}{2} & 2 & 0 \\[4pt]
    -\tfrac{1}{2} & 0 & 1
    \end{bmatrix},
    \]
    and note $\det(A) = -\tfrac{3}{4} \ne 0$, thus $C$ is non-singular.
\end{proof}

\begin{proof}
    We can write $y = \alpha x + \beta$ since $L$ is given. With
    $a=1,\ b=-3,\ c=2,\ d=-1,\ e=0,\ f=1$, substituting into $C$ gives
    \[F(x, \alpha x + \beta) 
    = a x^2 
    + b x\left(\alpha x + \beta\right)
    + c\left(\alpha x + \beta\right)^2 
    + d x 
    + e\left(\alpha x + \beta\right)
    + f = 0.\]
    Collecting gives a quadratic in $x$,
    \[
    A x^2 + B x + C = 0,
    \]
    where
    \[A = a + b\alpha + c\alpha^2, \quad B = b\beta + 2c\alpha\beta + d + e\alpha, \quad C = c\beta^2 + e\beta + f.\]
    We know $A \ne 0$ since $L \cap C$ consists of two points, so the quadratic
    has exactly two roots $x_0, x_1$ and factors as
    \[
    A x^2 + B x + C = A(x - x_0)(x - x_1).
    \]
    Expanding the right side gives $A x^2 - A(x_0 + x_1)x + A x_0 x_1$. Comparing the
    coefficient of $x$ with $B$ gives $-A(x_0 + x_1) = B$, thus
    \[
    x_1 = -\frac{B}{A} - x_0,
    \]
    and the $y$-coordinate is $y_1 = \alpha x_1 + \beta$. Thus the second point is
    \[
    P_1 = \left(-\frac{B}{A} - x_0,\ \ \alpha\left(-\frac{B}{A} - x_0\right) + \beta\right).
    \]
    This is clearly rational answering part (iii).
\end{proof}

\begin{tcolorbox}[title=Problem 3, breakable]
    Let $C$ be the conic given by the equation 
    \[F(x,y) = x^2 - 3xy + 2y^2 - x + 1 = 0.\]
    \begin{enumerate}
        \item Check that $C$ is non-singular (cf. the previous exercise).
        \item Let $L$ be the line $y = \alpha x + \beta$. Suppose that the intersection $L \cap C$ contains the point 
        $P_0 = (x_0, y_0)$. Assuming that $L \cap C$ consists of two distinct points, find the second point 
        of $L \cap C$ in terms of $\alpha, \beta, x_0, y_0$.
        \item If $L$ is a rational line and $P_0$ is a rational point (i.e., $\alpha, \beta, x_0, y_0 \in \mathbb{Q}$),
        prove that the second point of $L \cap C$ is also a rational point.
    \end{enumerate}
\end{tcolorbox}

\begin{tcolorbox}[title=Problem 4, breakable]
    Find all primitive integral right triangles 
    whose hyptenuse has length less than $30$.
\end{tcolorbox}

\begin{proof}
    We start by finding all relatively prime integers $>0$
    whose squares sum to $\le 30$ with $m < n$ and $m$ 
    and $n$ are not both odd.
    \begin{enumerate}
        \item $1^2 + 2^2 = 5 \le 30$
        \item $1^2 + 4^2 = 17 \le 30$
        \item $2^2 + 3^2 = 13 \le 30$
        \item $2^2 + 5^2 = 29 \le 30$
        \item $3^2 + 4^2 = 25 \le 30$
    \end{enumerate}
    For each  pair $(m,n)$, we get a triple
    $X = n^2 - m^2$, $Y = 2mn$, $Z = n^2 + m^2$ so
    \begin{enumerate}
        \item $(m,n) = (1,2)$ we have $(3, 4, 5)$
        \item $(m,n) = (1,4)$ we have $(15, 8, 17)$
        \item $(m,n) = (2,3)$ we have $(5, 12, 13)$
        \item $(m,n) = (2,5)$ we have $(21, 20, 29)$
        \item $(m,n) = (3,4)$ we have $(7, 24, 25)$
    \end{enumerate}
\end{proof}

\begin{tcolorbox}[title=Problem 5, breakable]
    Describe all rational points on the 
        circle $x^2 + y^2 = 2$.
\end{tcolorbox}

\begin{proof}
    Let $F(x, y) = x^2 + y^2 - 2 = 0$.
    Consider the line $L = y = -x$.
    Now we compute the line through $(1, 1)$ and $y = -x$.
    Now paramtrizing $y = -x$ gives the set of points of the form $(t, -t)$ where $t \in \mathbb{Q}$.
    So 
    \[y = \frac{-t - 1}{t - 1}(x - 1) + 1.\]
    Notice, since we take $x \ne 1$ the line is well defined.
    Then $F\left(x, \frac{-t - 1}{t - 1}(x - 1) + 1\right)$ we have 
    \[\frac{2(x - 1)((t^2 + 1)x - 2t)}{(t - 1)^2} = 0.\]
    Then, $x = 1$ corresponds to the known point $(1, 1)$ so we have 
    \[(t^2 + 1)x - 2t = 0,\]
    so $x = \frac{2t}{t^2 + 1}$ and $y = \frac{t^2 - 1}{t^2 + 1}$.
    So the set of all rational points on the circle $x^2 + y^2 = 2$ is 
    \[
    \left\{
    \left(\frac{2t}{t^2 + 1}, \frac{t^2 - 1}{t^2 + 1}\right)
    : t \in \mathbb{Q}
    \right\}
    \cup \{(1, 1)\}.
    \]
\end{proof}

\newpage
\begin{tcolorbox}[title=Problem 6, breakable]
    \begin{enumerate}
        \item     Let $a, b, c, d, e, f$ be non-zero real numbers.
    Use the substitution $t = \tan(\theta/2)$ to transform the integral 
    \[\int  \frac{a + b \cos \theta + c \sin \theta}{d + e \cos \theta + f \sin \theta} d\theta\]
    into the integral of a rational function of $t$.
        \item Evaluate the integral 
        \[
        \int  \frac{a + b \cos \theta + c \sin \theta}{1 + \cos \theta + \sin \theta} d \theta.
        \]
    \end{enumerate}
\end{tcolorbox}

\begin{proof}
    Substituting $\cos \theta = \frac{1 - t^2}{1 + t^2}$, $\sin \theta = \frac{2t}{1 + t^2}$, and $d\theta = \frac{2}{1 + t^2} \, dt$ we have
    \[
        \frac{a + b \cos \theta + c \sin \theta}{d + e \cos \theta + f \sin \theta} \, d\theta
        = \frac{a + b \frac{1 - t^2}{1 + t^2} + c \frac{2t}{1 + t^2}}{d + e \frac{1 - t^2}{1 + t^2} + f \frac{2t}{1 + t^2}} \cdot \frac{2}{1 + t^2} \, dt.
    \]
    Multiplying numerator and denominator of the fraction by $1 + t^2$ gives
    \[
        = \frac{a(1 + t^2) + b(1 - t^2) + 2ct}{d(1 + t^2) + e(1 - t^2) + 2ft} \cdot \frac{2}{1 + t^2} \, dt
        = \frac{2\bigl[(a - b)t^2 + 2ct + (a + b)\bigr]}{(1 + t^2)\bigl[(d - e)t^2 + 2ft + (d + e)\bigr]} \, dt.
    \]
\end{proof}

\begin{proof}
    Setting $d = e = f = 1$ so
    \[
        \int \frac{2\bigl[(a-b)t^2 + 2ct + (a+b)\bigr]}{(1+t^2)\bigl[2(t+1)\bigr]}\,dt
        = \int \frac{(a-b)t^2 + 2ct + (a+b)}{(t+1)(t^2+1)}\,dt
    \]
    \[
        = \frac{a+c}{2}\ln|t+1| + \frac{a-b}{4}\ln(t^2+1) + (b+c)\arctan t + K.
    \]
\end{proof}

\begin{tcolorbox}[title=Problem 7, breakable]
    For each of the following conics, either find a rational point or prove that there  
        are no rational points.
    \begin{enumerate}
        \item $x^2 + y^2 = 6$
        \item $3x^2 + 5y^2 = 4$
        \item $3x^2 + 6y^2 = 4$.
    \end{enumerate}
\end{tcolorbox}

\begin{proof}
    Suppose there exists $\frac{a}{b}, \frac{x}{y} \in \mathbb{Q}$
    such that $\left(\frac{a}{b}\right)^2 + \left(\frac{x}{y}\right)^2 = 6$.
    Then 
    \[\left(\frac{a}{b}\right)^2 + \left(\frac{x}{y}\right)^2 = 6
        \iff a^2 y^2 + x^2 b^2 = b^2 y^2 6.\]
\end{proof}